% ===========================================================================
%  E04 — Sala de aula: o professor abre a missão
% ===========================================================================
\chapter{E04 \textperiodcentered{} Sala de aula: o professor abre a missão}
\label{ep:e04}

\begin{fichaepico}
\begin{campos}
  \item[Objetivo] O professor conhece as missões, escolhe a da aula e a abre
    para a turma certa, no dia certo.
  \item[Para quem] Professor. A Coordenação consulta, sem mexer.
  \item[Valor] A missão acontece \emph{na aula}, com o professor presente. Se
    ela abrisse sozinha, o estudante jogaria em casa, sem o debate --- e o
    debate é onde a aprendizagem acontece.
  \item[Estado] \pronto{} --- falta a liberação com data marcada.
\end{campos}
\end{fichaepico}

\begin{tabularx}{\linewidth}{@{}L{16mm}X L{20mm}@{}}
\toprule
\headrow \thd{ID} & \thd{Feature} & \thd{Estado}\\
\midrule
\id{F04.1} & Minhas turmas & \pronto\\
\zebra \id{F04.2} & Liberar e encerrar missões na sala & \pronto\\
\id{F04.3} & Conhecer a trilha e a ficha da missão & \pronto\\
\zebra \id{F04.4} & Guia da aula & \pronto\\
\bottomrule
\end{tabularx}

% ---------------------------------------------------------------------------
\feature{F04.1}{Minhas turmas}{\pronto}

\begin{historia}{H04.1.1}{Começar pelo meu dia}{\pronto}
\hu{Professor}{ver no início os produtos que a minha rede contratou e as
minhas turmas}{ir direto ao que vou usar hoje}
\aceite
\begin{itemize}
  \item O Início (\enquote{Olá, \{nome\}}) mostra os cartões dos produtos:
    Gerador de Plano de Aula, Planejamento de Ensino, Entre Redes e Como você
    está?. O do bem-estar mostra a situação do dia (ex.: \enquote{Você já
    respondeu hoje \textperiodcentered{} 2 de 7}).
  \item Produto desligado para a rede não aparece --- nem como cartão
    desativado.
  \item \enquote{Minhas turmas} mostra as 4 primeiras, com \enquote{Ver
    todas}.
  \item \enquote{Continue de onde parou} mostra as atividades recentes do
    Instanooria.
\end{itemize}
\end{historia}

\begin{historia}{H04.1.2}{Ver só as turmas em que dou aula}{\pronto}
\hu{Professor}{ver as minhas turmas}{não me confundir com as turmas dos
colegas}
\aceite
\begin{itemize}
  \item \enquote{Minhas turmas} lista nome, série, turno, número de estudantes
    e período --- só das turmas em que estou matriculado.
  \item Sem turma, a tela orienta: \enquote{Quem faz o enturmamento é a
    coordenação da escola}.
  \item A turma de outro professor, aberta pelo endereço, não é encontrada.
\end{itemize}
\end{historia}

% ---------------------------------------------------------------------------
\feature{F04.2}{Liberar e encerrar missões na sala}{\pronto}

\begin{historia}{H04.2.1}{Liberar uma missão para a turma}{\pronto}
\hu{Professor}{liberar a missão para a minha turma no dia da aula}{que só essa
turma entre, e só quando eu estiver com ela}
\aceite
\begin{itemize}
  \item Na sala da turma, a aba Missões mostra a grade da série, com a
    etiqueta de cada missão: \enquote{aberta}, \enquote{fechada} ou
    \enquote{encerrada}.
  \item \enquote{Liberar} abre a missão, e a linha mostra \enquote{aberta por
    \{nome\} em \{data\}}.
  \item Só quem dá aula na turma consegue liberar.
  \item Sem liberação, o estudante vê o cartão com cadeado --- e digitar o
    endereço da missão também não abre.
  \item Sem nenhuma missão liberada, a sala avisa: \enquote{Nenhuma missão está
    aberta para esta turma}.
  \item Cada missão tem um único caminho a partir da sala: \enquote{Conduzir
    com esta turma}.
\end{itemize}
\end{historia}

\begin{historia}{H04.2.2}{Encerrar uma missão}{\pronto}
\hu{Professor}{encerrar a missão quando a sequência de aulas acabar}{a turma
seguir para a próxima}
\aceite
\begin{itemize}
  \item \enquote{Encerrar} fecha a entrada, mas não expulsa quem está jogando
    naquele momento.
  \item Nada do que a turma fez é apagado; a missão continua visível como
    \enquote{encerrada}.
\end{itemize}
\end{historia}

\begin{historia}{H04.2.3}{Programar a abertura e o fechamento}{\afazer}
\hu{Professor}{deixar a missão programada para abrir e fechar em datas
certas}{não depender de lembrar de liberar na hora da aula}
\aceiteproposta
\begin{itemize}
  \item Ao liberar, posso escolher \enquote{agora} ou uma janela (abre em
    \ldots{}, fecha em \ldots{}).
  \item A sala mostra a janela programada.
\end{itemize}
\end{historia}

\begin{historia}{H04.2.4}{Ver os estudantes da turma}{\pronto}
\hu{Professor}{ver quem está na turma e quem ainda não entrou}{ajudar quem está
com problema antes da aula}
\aceite
\begin{itemize}
  \item A aba Estudantes lista nome, usuário e matrícula, com as etiquetas
    \enquote{inativo} e \enquote{sem primeiro acesso}.
  \item O nome do estudante abre o perfil dele, com as habilidades da BNCC que
    ele já demonstrou (\id{H14.6.1}); ao lado, \enquote{Relatório da Vila do
    Lume} abre o relatório dele na Vila.
  \item O bloco \enquote{Também nesta turma} mostra os outros professores.
\end{itemize}
\end{historia}

% ---------------------------------------------------------------------------
\feature{F04.3}{Conhecer a trilha e a ficha da missão}{\pronto}

\begin{historia}{H04.3.1}{Consultar a trilha por ano}{\pronto}
\hu{Professor}{ver quais missões existem para cada ano}{planejar quando usar
cada uma}
\aceite
\begin{itemize}
  \item A trilha Entre Redes é organizada por etapa (Anos Iniciais, Anos
    Finais, Ensino Médio); o bloco \enquote{Para qualquer ano} vem primeiro.
  \item Anos com missão mostram \enquote{N de 4 unidades publicadas}; anos sem
    missão aparecem numa linha, \enquote{4 unidades em produção} --- sem
    título inventado.
  \item Faixa fixa: \enquote{Aqui você conhece as missões; quem abre para a
    turma é você, na sala}.
\end{itemize}
\end{historia}

\begin{historia}{H04.3.2}{Ler a ficha antes de levar para a turma}{\pronto}
\hu{Professor}{entender a missão antes de rodá-la}{conduzir sem estragar a
surpresa da revelação}
\aceite
\begin{itemize}
  \item A ficha tem \enquote{Como a missão funciona}, \enquote{Antes de rodar
    com a turma}, \enquote{O que o software NÃO conta ao estudante} e o
    roteiro.
  \item Botões \enquote{Abrir guia da aula} e \enquote{Abrir a aula e os
    resultados}.
  \item Vindo da sala, aparece a faixa \enquote{Conduzindo com a turma
    \{nome\}}.
  \item Unidade ainda não publicada: \enquote{\{missão\} ainda está em
    produção}.
\end{itemize}
\end{historia}

% ---------------------------------------------------------------------------
\feature{F04.4}{Guia da aula}{\pronto}

\begin{historia}{H04.4.1}{Ter um roteiro pronto para a aula}{\pronto}
\hu{Professor}{um roteiro de condução pronto para cada missão}{dar a aula sem
precisar prepará-la do zero}
\aceite
\begin{itemize}
  \item O guia tem nove seções fixas: objetivos e habilidade da BNCC
    Computação; preparo; o que não dizer; roteiro minuto a minuto; banco de
    perguntas; mediação; rubrica; adaptações; e o conteúdo real que a turma
    vê.
  \item Pode ser impresso.
  \item Existe para as nove missões publicadas.
\end{itemize}
\pendencia{o guia do unooGPT ainda não passou por revisão pedagógica em
sala.}
\end{historia}
